% 00_abstract.tex -- on the title page (not counted). Budget: ~200 words.
% problem -> idea (model both; sweep spill-over; breakdown points) -> C2 -> C1's teeth
%   -> biology (programmes, subtype dissociation) -> speed. C1 and C2 freeze after the
%   breakdown-gaps queue and the MintFlow refits (\pending below).
In imaging spatial transcriptomics, transcripts spilled over from neighbouring cells and a cell's response to its neighbours look alike in the counts: both make a cell resemble its niche. DISCELL models both. It splits each cell into an intrinsic state, kept free of the niche by a type-conditional adversary and graded by a held-out probe, and a response whose prior depends on the niche, and it mixes in a fraction $\kappa$ of the neighbours' expression. The counts bound $\kappa$ only from above, so we do not estimate it: we refit across $\kappa$ and report, for every finding, its breakdown point, the smallest spill-over fraction that explains it away. On four tissue sections and a serial section held out whole, no comparison method both carries less niche signal in its intrinsic state and keeps more of the cell's own state, on any section \pending{MintFlow refits}. The sweep has teeth: it breaks three findings, which are those spill-over could plausibly produce, while the others hold to four times the typical misassignment \pending{breakdown-gaps queue}. The response maps tissue territories with recognisable programmes, such as a stromal-remodelling programme around ovarian tumour, and held-out sublabels split as designed: tumour states are read from the intrinsic state, tumour- against stroma-associated locations from the response. One fit takes minutes, and the whole sweep costs less than one fit of the slowest comparison method.
